\section{Related Work}
\label{sec:related}

\textbf{Meme understanding.} The Hateful Memes challenge~\cite{kiela2020hateful} showed that unimodal models fail when
either modality alone is benign, and Memotion~\cite{sharma2020memotion} introduced the humour, sarcasm, offensiveness,
motivation and sentiment taxonomy that CMBan adopts. Most systems fuse a pretrained text encoder with a vision
encoder, often CLIP~\cite{radford2021clip}. Bangla resources include MemoSen~\cite{hossain2022memosen} (sentiment),
MUTE~\cite{hossain2022mute} and BHM~\cite{hossain2024bhm} (hate), BanglaAbuseMeme~\cite{das2023banglaabusememe}
(abuse, vulgarity, sarcasm, sentiment), MIMOSA~\cite{ahsan2024mimosa} (aggression) and CMBan~\cite{cmban2025}, which
annotates 2{,}641 cartoon memes along the Memotion dimensions. All evaluate on clean data only.

\textbf{Robustness.} ImageNet-C~\cite{hendrycks2019benchmarking} established severity-graded corruptions for vision,
CheckList~\cite{ribeiro2020checklist} behavioural testing for NLP, and GLUECoS~\cite{khanuja2020gluecos} showed that
multilingual encoders degrade on code-switched input. CLIP reads rendered text and is therefore susceptible to
typographic manipulation~\cite{goh2021multimodal}, which matters for memes, where the text \emph{is} part of the
image. For Bangla, the dominant real-world variation is Romanised writing, which has not been studied for memes.

\textbf{Explainability.} Integrated Gradients~\cite{sundararajan2017ig} and
Grad-CAM~\cite{selvaraju2017gradcam,chattopadhay2018gradcampp} are widely used to justify predictions. We use them
\emph{comparatively}, measuring how modality reliance and visual attention shift between clean and perturbed versions
of the same meme.
